\begin{lemma}
Let \(I\subseteq[n]\), put \(k=|I|\), and write
\[
\ell_R=|\pi_R(I)|,\qquad \ell_B=|\pi_B(I)|.
\]
Assume \(\varepsilon_1\ge0\),
\(\noderef{ParameterHierarchyEventualComparisons}
(\varepsilon,\varepsilon_1,\varepsilon_2,n_0)\), \(n>n_0\),
\[
k=\noderef{TwoBiteNaturalIndependenceScale}(1+\varepsilon,n),
\]
the rounded sample-parameter identities
\[
m=\noderef{TwoBiteNaturalReducedVertexCount}(n),
\qquad
p=\noderef{TwoBiteEdgeProbability}\!\left(\frac12,n\right),
\]
the event \(\noderef{FiberAndDegreeEvent}\), and the medium and small
closed-pair events.  Then the number of unclosed projected pairs satisfies
\[
f(\ell_R,\ell_B,k,p,n)-9\varepsilon_1k^2
\le |\noderef{TwoBiteUnclosedProjectedPairs}(\omega,I)|,
\]
where \(f=\noderef{ProjectionOpenPairFunction}\).
\end{lemma}
\begin{proof}
Let \(P_I=\noderef{TwoBiteProjectionPairSet}(\omega,I)\) and let
\(C_I\subseteq P_I\) be the set of projected pairs closed by some \(X_x(I)\),
as in \(\noderef{TwoBiteProjectionPairClosed}\).  Thus
\[
C_I=\{e\in P_I: \noderef{TwoBiteProjectionPairClosed}(\omega,I,e)\}.
\]
Let
\[
U_I=\noderef{TwoBiteUnclosedProjectedPairs}(\omega,I).
\]
By the definition of \(\noderef{TwoBiteUnclosedProjectedPairs}\), \(U_I\) is
the filter of \(P_I\) by the negation of
\(\noderef{TwoBiteProjectionPairClosed}\).  Equivalently,
\[
U_I=\{e\in P_I:e\notin C_I\}=P_I\setminus C_I.
\]
This identity is literal: for every tagged projected pair \(e\), membership in
\(U_I\) means \(e\in P_I\) together with the negation of the predicate defining
\(C_I\), while membership in \(C_I\) means \(e\in P_I\) together with that
predicate.  Hence \(U_I\cap C_I=\varnothing\), \(U_I\cup C_I=P_I\), and
\(C_I\subseteq P_I\).  The finite set \(P_I\) is therefore the disjoint union
\[
P_I=U_I\,\dot\cup\, C_I.
\]
Consequently
\[
|P_I|=|U_I|+|C_I|
\]
as natural cardinalities.  Since \(C_I\subseteq P_I\), the natural subtraction
\(|P_I|-|C_I|\) is the same natural number as \(|U_I|\); after casting this
cardinality identity to real numbers we may use
\[
|U_I|=|P_I|-|C_I|.
\]
By \(\noderef{TwoBiteProjectionPairSet}\), the projected-pair set is the
disjoint union, tagged by the two summands, of the red projected pairs and the
blue projected pairs.  Therefore
\[
|P_I|=\binom{\ell_R}{2}_{\mathbb R}
     +\binom{\ell_B}{2}_{\mathbb R}.
\]
Here \(\ell_R=|\pi_R(I)|\) and \(\ell_B=|\pi_B(I)|\) are cardinalities viewed
as real numbers, and the two summands are disjoint because the red and blue
projected pairs carry different tags.
Define the two structured huge opposite-color quantities by
\[
A_R=
\min\left\{\binom{k-\ell_R}{2}_{\mathbb R},
  \binom{pn}{2}_{\mathbb R}+\binom{k-\ell_R-pn}{2}_{\mathbb R}\right\},
\]
\[
A_B=
\min\left\{\binom{k-\ell_B}{2}_{\mathbb R},
  \binom{pn}{2}_{\mathbb R}+\binom{k-\ell_B-pn}{2}_{\mathbb R}\right\}.
\]

Now split the closed-pair set \(C_I\) according to both the witness class and
the color tag of the projected pair.  Define
\[
S_I(x)\Longleftrightarrow
\noderef{TwoBiteSmallClass}(\omega,\varepsilon,I,x),\quad
M_I(x)\Longleftrightarrow
\noderef{TwoBiteMediumClass}(\omega,\varepsilon,I,x),
\]
\[
L_I(x)\Longleftrightarrow
\noderef{TwoBiteLargeClass}(\omega,\varepsilon,I,x),\quad
H_I(x)\Longleftrightarrow
\noderef{TwoBiteHugeClass}(\omega,I,x).
\]
Take \(e\in C_I\).  By the definition of
\(\noderef{TwoBiteProjectionPairClosed}\), \(e\) may have several witnesses;
fix one such witness \(x\) for the following case split.  With this chosen
witness, either \(e=\operatorname{red}(q)\) and
\[
q\in\noderef{TwoBiteBasePairs}(\pi_R(X_x(I))),
\]
or \(e=\operatorname{blue}(q)\) and
\[
q\in\noderef{TwoBiteBasePairs}(\pi_B(X_x(I))).
\]
In either case \(X_x(I)=\noderef{TwoBiteX}(\omega,I,x)\) is a filter of \(I\),
so \(X_x(I)\subseteq I\) and \(|X_x(I)|\le |I|=k\).  Its cardinality is
nonnegative.  Now compare the real number \(|X_x(I)|\) successively with the
three cutoffs appearing in the definitions of \(S_I,M_I,L_I,H_I\).  If it is
at most the small cutoff, then \(x\in S_I\).  Otherwise, if it is at most the
large cutoff, then it is above the small cutoff and \(x\in M_I\).  Otherwise,
if it is at most the huge cutoff, then it is above the large cutoff and
\(x\in L_I\).  In the remaining case it is above the huge cutoff, and the
already proved inequality \(|X_x(I)|\le k=|I|\) supplies the upper bound
required in \(\noderef{TwoBiteHugeClass}\), so \(x\in H_I\).  Hence the chosen
witness lies in \(S_I\), \(M_I\), \(L_I\), or \(H_I\).

For the three non-huge classes set
\[
Q_{\mathcal A}
=\noderef{TwoBiteProjectedPairsClosedInClass}(\omega,I,\mathcal A),
\qquad
\mathcal A\in\{L_I,M_I,S_I\}.
\]
For huge witnesses, split by the color of the witness and the tag of the
projected pair.  Let \(H_R^{=}\) be the red-tagged projected pairs closed by
huge red witnesses, \(H_B^{=}\) the blue-tagged projected pairs closed by huge
blue witnesses, \(H_R^{\mathrm{opp}}\) the blue-tagged projected pairs closed
by huge red witnesses, and \(H_B^{\mathrm{opp}}\) the red-tagged projected
pairs closed by huge blue witnesses.  These are the four huge alternatives
after unfolding the two `Sum` tags in
\(\noderef{TwoBiteProjectionPairClosed}\).

For huge witnesses the base vertex \(x\) is either \(x=\operatorname{inl}(r)\)
or \(x=\operatorname{inr}(b)\).  In the first case the same-color projection is
\(\pi_R\) and the opposite-color projection is \(\pi_B\); in the second case
the same-color projection is \(\pi_B\) and the opposite-color projection is
\(\pi_R\).  Equivalently, the two huge filters used below are exactly the
Lean predicates
\[
\mathsf{hugeRed}(x)\Longleftrightarrow H_I(x)\land
  \exists r,\ x=\operatorname{inl}(r),
\qquad
\mathsf{hugeBlue}(x)\Longleftrightarrow H_I(x)\land
  \exists b,\ x=\operatorname{inr}(b).
\]
The same-color huge pieces are generated by
\((\mathsf{hugeRed},\pi_R)\) and \((\mathsf{hugeBlue},\pi_B)\); the
opposite-color huge pieces are generated by
\((\mathsf{hugeRed},\pi_B)\) and \((\mathsf{hugeBlue},\pi_R)\).  Therefore
the preceding classification gives the set inclusion
\[
C_I\subseteq
Q_{L_I}\cup Q_{M_I}\cup Q_{S_I}\cup
H_R^{=}\cup H_B^{=}\cup
H_R^{\mathrm{opp}}\cup H_B^{\mathrm{opp}}. \tag{1}
\]
Indeed, the chosen witness for \(e\) is either large, medium, small, or huge.
In the first three cases \(e\) lies in the corresponding \(Q_{\mathcal A}\).
In the huge case, the witness is red or blue and the projected pair is
red-tagged or blue-tagged, giving exactly one of the four displayed huge
sets.  Since one arbitrary witness was chosen only to prove membership in the
right-hand side, (1) is a containment and does not assert that the displayed
sets are disjoint.  Define
\[
C_{\mathrm{err}}=
Q_{L_I}\cup Q_{M_I}\cup Q_{S_I}\cup H_R^{=}\cup H_B^{=},
\qquad
C_R^{\mathrm{opp}}=H_R^{\mathrm{opp}},\qquad
C_B^{\mathrm{opp}}=H_B^{\mathrm{opp}}.
\]
Then (1) implies
\[
C_I\subseteq
C_{\mathrm{err}}\cup C_R^{\mathrm{opp}}\cup C_B^{\mathrm{opp}}.
\]
Taking cardinalities and applying the finite union bound, without any
disjointness assumption, gives
\[
|C_I|\le |C_{\mathrm{err}}|
       +|C_R^{\mathrm{opp}}|+|C_B^{\mathrm{opp}}|. \tag{2}
\]

We first bound \(C_{\mathrm{err}}\).  Introduce the numerical contribution
\[
\begin{aligned}
D_{\mathrm{err}}
&=|\noderef{TwoBiteClosedPairsInClass}(I,L_I)|
  +|\noderef{TwoBiteClosedPairsInClass}(I,M_I)|
  +|\noderef{TwoBiteClosedPairsInClass}(I,S_I)|\\
&\quad+
\sum_{x\in H_I\cap V_R}
  \binom{|\pi_R(X_x(I))|}{2}_{\mathbb R}
+
\sum_{x\in H_I\cap V_B}
  \binom{|\pi_B(X_x(I))|}{2}_{\mathbb R}.
\end{aligned}
\]
Equivalently, write
\[
A_L=|\noderef{TwoBiteClosedPairsInClass}(\omega,I,L_I)|,\quad
A_M=|\noderef{TwoBiteClosedPairsInClass}(\omega,I,M_I)|,\quad
A_S=|\noderef{TwoBiteClosedPairsInClass}(\omega,I,S_I)|
\]
and let \(A_H^{=}\) be the displayed same-color huge projected sum; then
\[
D_{\mathrm{err}}=A_L+A_M+A_S+A_H^{=}.
\]
Here \(A_H^{=}\) is exactly the same expression that appears in the
same-color huge input of \(\noderef{ClosedPairsControlled}\): if
\[
\mathrm{hugeRed}(x)\Longleftrightarrow
\noderef{TwoBiteHugeClass}(\omega,I,x)\land x\in V_R,
\qquad
\mathrm{hugeBlue}(x)\Longleftrightarrow
\noderef{TwoBiteHugeClass}(\omega,I,x)\land x\in V_B,
\]
then, after unfolding \(\noderef{TwoBiteProjectedPairSum}\),
\[
A_H^{=}
=\noderef{TwoBiteProjectedPairSum}(\omega,I,\mathrm{hugeRed},\pi_R)
 +\noderef{TwoBiteProjectedPairSum}(\omega,I,\mathrm{hugeBlue},\pi_B).
\]
The first three summands count final-vertex pairs, not tagged projected pairs,
so they cannot be used with multiplicity one.  For each class
\(\mathcal A\in\{L_I,M_I,S_I\}\), the already-defined set
\(Q_{\mathcal A}\) is the projected-pair set generated by witnesses in
\(\mathcal A\).
The part of \(C_{\mathrm{err}}\) with a non-huge witness is contained in
\[
Q_{L_I}\cup Q_{M_I}\cup Q_{S_I}.
\]
For each of the three classes, apply
\(\noderef{ProjectedPairsInClassMultiplicityBound}\) with the same
\(\omega\), the same \(I\), and the displayed class \(\mathcal A\).  It gives
\[
|Q_{\mathcal A}|
\le
2|\noderef{TwoBiteClosedPairsInClass}(\omega,I,\mathcal A)|.
\]
This is the required conversion from final-pair counts to projected-pair
counts, and it is made before invoking \(\noderef{ClosedPairsControlled}\).
The non-huge part of \(C_{\mathrm{err}}\) is contained in
\[
Q_{L_I}\cup Q_{M_I}\cup Q_{S_I},
\]
so the finite union bound gives
\[
\left|Q_{L_I}\cup Q_{M_I}\cup Q_{S_I}\right|
\le |Q_{L_I}|+|Q_{M_I}|+|Q_{S_I}|.
\]
Substituting the three multiplicity inequalities, one for each of
\(L_I,M_I,S_I\), gives the explicit projected-pair estimate
\[
\left|Q_{L_I}\cup Q_{M_I}\cup Q_{S_I}\right|
\le 2A_L+2A_M+2A_S. \tag{3}
\]
The same-color huge part is already bounded by a projected-pair sum.  For a
fixed huge red witness \(x\), the red tagged pairs it closes are exactly the
red tags of the elements of
\(\noderef{TwoBiteBasePairs}(\pi_R(X_x(I)))\), so their number is
\(\binom{|\pi_R(X_x(I))|}{2}_{\mathbb R}\).  The union over all huge red
witnesses therefore has cardinality at most the corresponding sum over
\(H_I\cap V_R\).  The same argument for huge blue witnesses gives the sum over
\(H_I\cap V_B\).  Thus
\[
|H_R^{=} \cup H_B^{=}|\le A_H^{=}. \tag{4}
\]
Combining (3) and (4), and taking a union bound over the five possible sources
of \(C_{\mathrm{err}}\), with overlaps only decreasing the union cardinality,
gives
\[
|C_{\mathrm{err}}|
\le
2A_L+2A_M+2A_S+A_H^{=}.
\]
Since \(A_L,A_M,A_S,A_H^{=}\) are all nonnegative, the last display is at most
\[
2(A_L+A_M+A_S+A_H^{=})=2D_{\mathrm{err}}.
\]
Hence
\[
|C_{\mathrm{err}}|\le 2D_{\mathrm{err}}.
\]

We now verify the inputs to \(\noderef{ClosedPairsControlled}\) with the same
parameters used in the statement.  The set-size input is the equality
\(|I|=k\), where
\[
k=\noderef{TwoBiteNaturalIndependenceScale}(1+\varepsilon,n).
\]
The degree input is the assumed \(\noderef{FiberAndDegreeEvent}(\omega,\varepsilon_2)\).
The medium and small inputs are exactly the two assumed events
\(\noderef{TwoBiteMediumClosedPairsEvent}\) and
\(\noderef{TwoBiteSmallClosedPairsEvent}\), both at scale \(k=|I|\) and at the
same \(\omega,\varepsilon,\varepsilon_1\).  Unfolding those two event
definitions gives the two \(\noderef{ClosedPairsBound}\) hypotheses for
\(A_M\) and \(A_S\) that \(\noderef{ClosedPairsControlled}\) requires.

The large input required by \(\noderef{ClosedPairsControlled}\) is supplied by
\(\noderef{LargeClosedPairsBound}\), applied to the same
\(\omega,I,\varepsilon,\varepsilon_1,\varepsilon_2,n_0\).  Its hypotheses are
precisely the eventual-comparisons package, \(n>n_0\), the rounded scale
identity \(|I|=\noderef{TwoBiteNaturalIndependenceScale}(1+\varepsilon,n)\),
and the same fiber-and-degree event.  Hence it gives
\[
\noderef{ClosedPairsBound}
\left(|\noderef{TwoBiteClosedPairsInClass}(\omega,I,L_I)|,\varepsilon_1,k\right).
\]
The same-color huge input is the third conjunct of
\(\noderef{HugeProjectionClosedPairsBound}\), applied to the same
\(\omega,I,\varepsilon,\varepsilon_1,\varepsilon_2,n_0\), using the assumed
nonnegativity \(\varepsilon_1\ge0\), the same eventual-comparisons package,
the same strict inequality \(n>n_0\), the rounded identity
\[
k=|I|=\noderef{TwoBiteNaturalIndependenceScale}(1+\varepsilon,n),
\]
the same rounded sample-parameter identities
\[
m=\noderef{TwoBiteNaturalReducedVertexCount}(n),\qquad
p=\noderef{TwoBiteEdgeProbability}\!\left(\frac12,n\right),
\]
the equality \(|I|=k\), and the same \(\noderef{FiberAndDegreeEvent}\).  It
gives
\[
\noderef{ClosedPairsBound}(A_H^{=},\varepsilon_1,k).
\]
Thus all hypotheses of \(\noderef{ClosedPairsControlled}\) are satisfied.
Its conclusion, with the four grouped summands in exactly the order
\((L_I,M_I,S_I,H_I\text{ same-color})\), is
\[
\noderef{ClosedPairsBound}(D_{\mathrm{err}},4\varepsilon_1,k).
\]
By the definition of \(\noderef{ClosedPairsBound}\), this is exactly
\[
D_{\mathrm{err}}\le 4\varepsilon_1k^2.
\]
This invocation comes after the multiplicity step above: the child controls
the final-pair and same-color huge numerical sum \(D_{\mathrm{err}}\), while
\(\noderef{ProjectedPairsInClassMultiplicityBound}\) has already converted
the non-huge projected-pair pieces into the factor-two contribution.  Together
with \(|C_{\mathrm{err}}|\le 2D_{\mathrm{err}}\), this proves the
corrected error estimate
\[
|C_{\mathrm{err}}|\le 8\varepsilon_1k^2.
\]

The remaining two pieces are controlled by the huge opposite-color projected
sums.  Let
\[
S_R^{\mathrm{opp}}=
\sum_{x\in H_I\cap V_R}
  \binom{|\pi_B(X_x(I))|}{2}_{\mathbb R},
\qquad
S_B^{\mathrm{opp}}=
\sum_{x\in H_I\cap V_B}
  \binom{|\pi_R(X_x(I))|}{2}_{\mathbb R}.
\]
For a fixed huge red witness, the opposite-color projected pairs it closes are
the blue tags of
\(\noderef{TwoBiteBasePairs}(\pi_B(X_x(I)))\), so their number is
\(\binom{|\pi_B(X_x(I))|}{2}_{\mathbb R}\).  Hence the union over all huge red
witnesses, namely \(C_R^{\mathrm{opp}}\), has cardinality at most
\(S_R^{\mathrm{opp}}\).  Similarly
\(|C_B^{\mathrm{opp}}|\le S_B^{\mathrm{opp}}\).  These are union bounds over
witnesses, not disjoint decompositions; multiple witnesses for one projected
pair only make the sums larger.
Apply \(\noderef{HugeProjectionClosedPairsBound}\) a second time with the same
\(\omega,I,\varepsilon,\varepsilon_1,\varepsilon_2,n_0\), the same
nonnegativity, eventual-comparisons package, strict inequality \(n>n_0\),
rounded identity \(k=\noderef{TwoBiteNaturalIndependenceScale}(1+\varepsilon,n)\),
the same rounded sample-parameter identities
\[
m=\noderef{TwoBiteNaturalReducedVertexCount}(n),\qquad
p=\noderef{TwoBiteEdgeProbability}\!\left(\frac12,n\right),
\]
the same equality \(|I|=k\), and the same \(\noderef{FiberAndDegreeEvent}\).
Its fourth conjunct is the huge-red-through-blue estimate: its left side uses
the huge-red predicate but the opposite projection \(\pi_B\), while its
deficit parameter is the size of the same-color red projection of \(I\).  Since
\[
\ell_R=|\pi_R(I)|,
\]
this conjunct is exactly
\[
S_R^{\mathrm{opp}}\le(1+\varepsilon_1)A_R.
\]
Its fifth conjunct is the huge-blue-through-red estimate: its left side uses
the huge-blue predicate but the opposite projection \(\pi_R\), while its
deficit parameter is the size of the same-color blue projection of \(I\).
Because
\[
\ell_B=|\pi_B(I)|,
\]
it is exactly
\[
S_B^{\mathrm{opp}}\le(1+\varepsilon_1)A_B.
\]
Hence
\[
|C_R^{\mathrm{opp}}|\le (1+\varepsilon_1)A_R,
\qquad
|C_B^{\mathrm{opp}}|\le (1+\varepsilon_1)A_B.
\]
Combining these two bounds with the union bound for \(C_I\) gives
\[
|C_I|
\le
8\varepsilon_1k^2+(1+\varepsilon_1)A_R+(1+\varepsilon_1)A_B
=A_R+A_B+8\varepsilon_1k^2+\varepsilon_1(A_R+A_B).
\]
We now bound the final \(\varepsilon_1\)-multiple.  Since
\(\pi_R(I)\) and \(\pi_B(I)\) are images of the \(k\)-set \(I\), we have
\[
\ell_R\le k,\qquad \ell_B\le k.
\]
Let \(d_R=k-|\pi_R(I)|\) and \(d_B=k-|\pi_B(I)|\) as natural-cardinality
deficits; the displayed inequalities show \(0\le d_R,d_B\le k\), and after
casting to real numbers \(d_R=k-\ell_R\) and \(d_B=k-\ell_B\).  Choosing
inclusions of a \(d_R\)-element set and a \(d_B\)-element set into a fixed
\(k\)-element set gives injections on unordered pairs, so
\[
\binom{d_R}{2}_{\mathbb R}\le \binom{k}{2}_{\mathbb R},
\qquad
\binom{d_B}{2}_{\mathbb R}\le \binom{k}{2}_{\mathbb R}.
\]
The structured quantities \(A_R\) and \(A_B\) are minima.  A minimum is no
larger than its first argument, so
\[
A_R\le \binom{d_R}{2}_{\mathbb R},\qquad
A_B\le \binom{d_B}{2}_{\mathbb R}.
\]
Combining the last two displays gives
\[
A_R\le \binom{k}{2}_{\mathbb R}\le \frac{k^2}{2},
\qquad
A_B\le \binom{k}{2}_{\mathbb R}\le \frac{k^2}{2}.
\]
The last inequalities use
\(\binom{k}{2}_{\mathbb R}=k(k-1)/2\le k^2/2\) for the nonnegative
cardinality \(k\).  Adding the two bounds gives \(A_R+A_B\le k^2\), and
because \(\varepsilon_1\ge0\),
\[
\varepsilon_1(A_R+A_B)\le \varepsilon_1k^2.
\]
Therefore
\[
|C_I|\le A_R+A_B+9\varepsilon_1k^2.
\]

Finally, by the definition of \(f=\noderef{ProjectionOpenPairFunction}\),
\[
f(\ell_R,\ell_B,k,p,n)
=\binom{\ell_R}{2}_{\mathbb R}+\binom{\ell_B}{2}_{\mathbb R}-A_R-A_B
=|P_I|-A_R-A_B.
\]
Combining this identity with the previous bound gives
\[
|U_I|
=|P_I|-|C_I|
\ge f(\ell_R,\ell_B,k,p,n)-9\varepsilon_1k^2,
\]
as required.
\end{proof}
